
\lesson{11/11}
\textbf{Note} For algebras of Hermitian matrices, the anticommutator
is anti-hermitian:
\[
    \left(\left[X_a,X_b\right]\right)^\dagger =[...]=-\left[X_a,X_b\right]
\]
\subsection{Compact, simple Lie Groups}
\textbf{Properties:}
\begin{itemize}
    \item All finite dimensional representations are \textit{equivalent} to
        unitary representations;
    \item All reducible unitary representations are \textit{completely reducible}.
\end{itemize}

For each group of this kind there are two important representations:
\begin{enumerate}
    \item The \textbf{defining representation} (e.g. the matrix definition of
        $SU(n)$);
    \item The \textbf{adjoint representation}: given an element of the \textit{
        algebra} of the group $X_a$, its adjoint representation is
        \[
            \left[D^\text{adj}\left(X_a\right)\right]_{bc}= -if_{abc} 
        \]
        which means this representation has \textit{the same dimension as the
        Lie Algebra}.
\end{enumerate}
\textbf{Example:} $\mathbf{SU(2)}$\\
The generators of the algebra are the Pauli matrices, imposing
\[
    g=\mathbb{I}+itA+O(t^2)
\]
yields
\[
    J_i=\frac{1}{2}\sigma_i.
\]
The commutation relations are
\[
    \left[J_a,J_b\right]=i\varepsilon_{abc}J_c
\]
So the structure constants are $f_{abc}=i\varepsilon_{abc}$ and the
adjoint representation of the basis of the algebra is then
\[
    J_1^\text{adj}=\begin{pmatrix}
        0&0&0\\
        0&0&-i\\
        0&i&0
    \end{pmatrix},\;J_2^\text{adj}=\begin{pmatrix}
        0&0&i\\
        0&0&0\\
        -i&0&0
    \end{pmatrix},\;J_3^\text{adj}=\begin{pmatrix}
        0&-i&0\\
        i&0&0\\
        0&0&0
    \end{pmatrix} 
\]

How can irreps be characterized?

\subsubsection{Dynkin Index}
Given a basis $\left\{X_a\right\}$ of the Lie algebra, define
\[
    M\in\text{mat}(\dim\mathfrak{g}\times \dim\mathfrak{g})\text{ s.t. }M_{ab}=
    \tr(X_aX_b)
\]
$M$ is clearly symmetric, so it can be diagonalized
\[
    M_{ab} = \lambda_d\delta_{ab}
\]
the number $\lambda_d$ is the Dynkin index [??unclear??]\\
something with $\otimes\oplus$.
\subsubsection{Second-order Casimir operator}
Consider the operator
\[
    \rho^{(2)} = \sum_{a=1}^{\dim\mathfrak{g}}X_aX_b
\]
$\forall b,\; \left[\rho,X_b\right]=0$, so the operator is proportional to the
identity:
\[
    \rho^{(2)}=C_2(D)\mathbb{I} 
\]
The trace of this operator is simply
\[
    \Tr\rho^{\left(2\right)} = \sum_{a}\Tr(X_aX_a) =\dim\mathfrak{g}\cdot \lambda_d
\]
but also
\[
    \Tr\rho ^{(2)} = C_2(D)\dim(D)
\]
so
\[
    \lambda_D=\frac{\dim(D)C_2(D)}{\dim\mathfrak{g}}
\]
more things missing...
\subsection{Roots and weights}
Consider a Lie algebra $\mathfrak{g}$ of dimension $N$ and \textit{the set
of its generators that commute with each other}
\[
    \left\{H_1,...,H_m\right\}\subset\mathfrak{g}
\]
with $1\leq m\leq N$: these generators are called \textit{Cartan generators} and
$m$ is called the \textit{rank} of the Lie Algebra. These generators can be
rescaled so that
\[
    \Tr(H_iH_j)=K_D\delta_{ij}
\]
where $k_D$ depends only on the representation $D$. Consider now the adjoint
representation: its dimension is $N$ and its generators are $\left\{X_a\right\}$.
Associate now to every generator $X_a$ a vector in the space the generators act on
\[
    X_a\to\ket{X_a}
\]
and construct the scalar product
\[
    \bra{X_a}\ket{X_b}:=\Tr(X_a^\dagger X_b).
\]
One can choose the vectors $\ket{X_a}$ so that
\[
    X_a\ket{X_b} = if_{abc}\ket{X_c}:=\ket{\left[X_a,X_b\right]}
\]
this notation is consistent, since mapping the combination of generators
$if_{abc}X_c$ to a vector yields the same result.

Doing the same for the Cartan generators $H_i$ on the same vector space yields
\[
    H_a\ket{H_b}=\ket{\left[H_a,H_b\right]}=0
\]
Consider the common eigenvector of all $H_i$:
\[
    H_i\ket{E_{\alpha}}=\alpha_i\ket{E_{\alpha}},\;\alpha \in\mathbb{R}^m 
\]
Expressing the eigenvector as a linear combination of $\ket{X_a}$ means to
express $E_\alpha$ as a linear combination of $X_a$! This means
\[
    \left[H_i,E_{\alpha}\right]=\alpha_iE_\alpha 
\]
The eigenvalue equation holds for the adjoint representation, while this second 
relation \textbf{holds in general}.\\
The vectors $\ket{E_\alpha}$ are called \textit{roots} of the adjoint
representation.

Consider now a representation $D$: the eigenvectors of the Cartan generators in
this rep. can be denoted as
\[
    \ket{\mu,D}\to H_i^{(D)}\ket{\mu,D}=\mu_i\ket{\mu,D}
\]
where $\mu\in\mathbb{R}^m$; consider now the quantity (always in $D$)
\[
    H_iE_\alpha\ket{\mu,D} = (H_iE_\alpha-E_\alpha H_i+E_\alpha H_i)\ket{\mu,D}=
    \left[H_i,E_\alpha\right]\ket{\mu,D} + \mu_iE_\alpha\ket{\mu,D}
\]
The result of the commutator is independent of $D$, so
\[
    H_iE_\alpha\ket{\mu,D} = (\alpha_i+\mu_i)E_\alpha \ket{\mu,D}
\]
so applying $E^{(D)}_\alpha$ to an eigenvector of $H_i^{(D)}$ yields another
eigenvector of $H_i$:
\[
    E_\alpha \ket{\mu,D} = N_{\alpha ,\mu}\ket{\alpha +\mu,D}
\]
What about the quantity
\[
    H_iE_\alpha \ket{E_{-\alpha }}
\]
in the adjoint representation?
\[
    H_iE_\alpha \ket{E_{-\alpha}}=\left[H_i,E_\alpha\right]\ket{E_{-\alpha }}+
    E_{\alpha }H_i\ket{E_{-\alpha }} =  0
\]
So $E_\alpha \ket{E_{-\alpha}} = \sum_{i=1}^{m} \alpha _iH_i$. [???]

In a generic representation $D$,
\[
    \bra{\mu,D}\left[E_\alpha,E_{-\alpha}\right]\ket{\mu,D}=\sum_{i}^{}\alpha_i
    \bra{\mu,D}H_i\ket{\mu,D} = \alpha\cdot \mu
\]
{\large\textbf{[...other holes...guarda dal libro, manca un'ora di lezione qui]}
}\\
`
The fundamental result is
\[
    \frac{a\cdot\mu}{\alpha^2}=\frac{M}{2}
\]
where $M=q-p\in\mathbb{Z}$ 
